\section{The Exclamation Mark Nobody Typed}
\label{sec:ch14-the-exclamation-mark-nobody-typed}

\index{node@\code{node}!the plural field's wrapper|(}%
Every non-null field so far has been one somebody wrote. There is one in this
graph that nobody did.

\begin{graphqlsdl}
type Query {
  "Fetches an object given its ID."
  node("ID of the object." id: ID!): Node @shareable
  "Lookup nodes by a list of IDs."
  nodes("The list of node IDs." ids: [ID!]!): [Node]! @shareable
}
\end{graphqlsdl}

\index{AddGlobalObjectIdentification@\code{AddGlobalObjectIdentification}!generates two different nullabilities}%
Both fields come from \code{AddGlobalObjectIdentification()}, which
chapter~\ref{ch:schema-design} called once and has not been thought about
since. The singular field is nullable and the plural one is a non-null list of
nullable items. Nobody in this book chose either, and the difference between
them turns out to be the difference between losing a field and losing the
response.

\subsection{The Item Type Does Not Save You}

\index{node@\code{node}!nullable items promise nothing}%
\code{[Node]} inside \code{[Node]!} reads like a design that has already thought
about partial failure: the list must be there, and any entry in it may be null,
so one unreachable object should cost one entry. Stop the Speakers service and
ask for one session and one speaker.

\begin{minted}{graphql}
query N($i: [ID!]!) {
  nodes(ids: $i) {
    __typename
    ... on Speaker { name }
    ... on Session { title }
  }
}
\end{minted}

The variable holds one session id and one speaker id, in that order:

\begin{minted}{json}
{"i": ["U2Vzc2lvbjox", "U3BlYWtlcjox"]}
\end{minted}

The router answers:

\begin{minted}[fontsize=\footnotesize,breakanywhere]{json}
{"errors":[{"message":"Failed to fetch from Subgraph 'speakers'."}],"data":null}
\end{minted}

\index{node@\code{node}!no path, no index, no partial list}%
There is no list with a hole in it, and the error carries neither an index nor
a path. The session in that request is answered by a service that is running,
and asking for it alone proves it:

\begin{minted}[fontsize=\footnotesize,breakanywhere]{json}
{"data":{"nodes":[{"__typename":"Session","title":"Schemas That Outlive Their Authors"}]}}
\end{minted}

\index{router!fails the field rather than the entry}%
The reason the mixed list does not degrade is that the router does not resolve
\code{nodes} entry by entry. It plans one batched step per owning subgraph, and
when a step has nowhere to go it fails the field rather than the entries under
it. A failure attributed to the whole field lands on the field's own type, that
type is \code{[Node]!}, and the walk from
section~\ref{sec:ch14-how-far-a-null-climbs} takes it to \code{data}.

\index{node@\code{node}!the singular field contains the same failure}%
Now the singular field, same dead service, same speaker:

\begin{minted}[fontsize=\footnotesize,breakanywhere]{json}
{"errors":[{"message":"Failed to fetch from Subgraph 'speakers'."}],"data":{"node":null}}
\end{minted}

\index{blast radius!decided by the wrapper, not the failure}%
Identical failure, identical error text, and the damage stops at the field
because \code{node} is \code{Node} rather than \code{Node!}. Two generated
fields, one line apart in the same schema, and the plural one is the one that
can empty a response.

\subsection{Reaching a Field You Did Not Declare}

\index{TypeInterceptor@\code{TypeInterceptor}!the second thing it is needed for}%
There is no signature to put a question mark on. \code{node} and \code{nodes}
are generated and declared in no class of ours, which is the same wall
chapter~\ref{ch:composition} hit when it needed to mark those two fields
shareable and answered with a type interceptor. The field's type is settable on
the configuration the interceptor is handed, so the same hook works for this:

\begin{minted}{csharp}
using HotChocolate.Configuration;
using HotChocolate.Types.Descriptors;
using HotChocolate.Types.Descriptors.Configurations;

namespace Sessions;

/// <summary>
/// Makes the generated nodes field nullable, so that a subgraph it
/// routes to being down costs the field rather than the response.
/// </summary>
// AddGlobalObjectIdentification generates node and nodes, and it
// generates them with different nullability: node is Node, which may
// hold a null, and nodes is [Node]!, which may not. The item type
// inside that list is nullable and it does not help, because the
// router resolves nodes as one batched step per owning subgraph and
// fails the whole field when one of them is unreachable. A non-null
// list with nowhere to put the failure propagates to its parent, and
// the parent of a root field is the response.
//
// Neither field is declared in any class here, so there is no
// signature to put a question mark on. This is the same problem
// ShareNodeFields beside it solves for a directive, and the same hook
// solves it: the field's Type is settable on the configuration.
internal sealed class NullableNodesField : TypeInterceptor
{
    public override void OnBeforeCompleteType(
        ITypeCompletionContext context,
        TypeSystemConfiguration configuration)
    {
        if (configuration is not ObjectTypeConfiguration type
            || type.Name != "Query")
        {
            return;
        }

        foreach (var field in type.Fields)
        {
            if (field.Name is "nodes")
            {
                field.Type = TypeReference.Parse(
                    "[Node]", TypeContext.Output);
            }
        }
    }
}
\end{minted}

\index{TryAddTypeInterceptor@\code{TryAddTypeInterceptor}!registers the second one}%
Registering it is one more line on the builder, below the
chapter~\ref{ch:composition} one. Here is the Sessions service's
\code{Program.cs} again, complete, with this chapter's line marked:

\begin{minted}[highlightlines={37}]{csharp}
using ChilliCream.Nitro.App;
using HotChocolate.AspNetCore;
using Microsoft.EntityFrameworkCore;
using Sessions;

var builder = WebApplication.CreateBuilder(args);

// EF Core reports every command through the logging pipeline as well,
// at Information and with a duration attached. StatementLog below is
// this service's instrument, so silence the other one rather than have
// every statement printed twice.
builder.Logging.AddFilter(
    "Microsoft.EntityFrameworkCore.Database.Command",
    LogLevel.Warning);

builder.Services.AddDbContextFactory<ConferenceContext>(options =>
{
    options.UseSqlite("Data Source=conference.db");

    if (builder.Environment.IsDevelopment())
    {
        options.AddInterceptors(new StatementLog());
    }
});

builder
    .AddGraphQL()
    .AddSessionsTypes()
    .RegisterDbContextFactory<ConferenceContext>()
    .AddQueryContext()
    .AddPagingArguments()
    .AddGlobalObjectIdentification()
    .AddMutationConventions()
    .AddApolloFederation()
    .ModifyCostOptions(options => options.ApplyCostDefaults = false)
    .TryAddTypeInterceptor<ShareNodeFields>()
    .TryAddTypeInterceptor<NullableNodesField>();

var app = builder.Build();

using (var scope = app.Services.CreateScope())
{
    var factory = scope.ServiceProvider
        .GetRequiredService<IDbContextFactory<ConferenceContext>>();
    using var db = factory.CreateDbContext();
    db.Database.EnsureCreated();
}

app.MapGraphQL()
    .WithOptions(options =>
    {
        // Serve the Nitro build that shipped in the package. The
        // default, ServeMode.Latest, downloads the current one
        // from ChilliCream's CDN instead.
        options.Tool.ServeMode = ServeMode.Embedded;
    });

app.RunWithGraphQLCommands(args);
\end{minted}

\index{@shareable@\code{@shareable}!two subgraphs must agree on the type}%
That is half the job, and I expected the composer to insist on the other half.
\code{node} and \code{nodes} are shareable, which says two subgraphs can
resolve the same field, and a field two subgraphs declare with different types
does not look like one field. Soften it in Sessions, leave Speakers at
\code{[Node]!}, and compose.

\index{@shareable@\code{@shareable}!the composer takes the permissive type}%
It composes. Exit zero, nothing printed, and the schema the router serves
carries \code{[Node]}: handed two types for one shareable field, this composer
takes the more permissive of them rather than refusing the pair. That is the
third time in this book that this composer has accepted something a reader
would expect it to stop, after the \code{@cost} weights of
chapter~\ref{ch:composition} and the \code{@provides} preconditions of
chapter~\ref{ch:second-third-subgraph}.

\index{@shareable@\code{@shareable}!why both services change anyway}%
Both services change anyway. A graph that is correct because a composer was
lenient is a graph that breaks on the release where it stops being, and the
half-softened version leaves the Speakers subgraph telling anyone who reads its
schema a different story from the one the supergraph tells. So Speakers gets
its own copy, exactly as it got its own \code{ShareNodeFields} in
chapter~\ref{ch:second-third-subgraph}:

\begin{minted}{csharp}
using HotChocolate.Configuration;
using HotChocolate.Types.Descriptors;
using HotChocolate.Types.Descriptors.Configurations;

namespace Speakers;

/// <summary>
/// Makes the generated nodes field nullable, so that a subgraph it
/// routes to being down costs the field rather than the response.
/// Declared again in this service rather than shared with the
/// Sessions one, for the same reason ShareNodeFields is: both
/// subgraphs export this field, so both have to soften it, and the
/// two must agree or the composed field has two types.
/// </summary>
// AddGlobalObjectIdentification generates node and nodes, and it
// generates them with different nullability: node is Node, which may
// hold a null, and nodes is [Node]!, which may not. The item type
// inside that list is nullable and it does not help, because the
// router resolves nodes as one batched step per owning subgraph and
// fails the whole field when one of them is unreachable. A non-null
// list with nowhere to put the failure propagates to its parent, and
// the parent of a root field is the response.
//
// Neither field is declared in any class here, so there is no
// signature to put a question mark on. This is the same problem
// ShareNodeFields beside it solves for a directive, and the same hook
// solves it: the field's Type is settable on the configuration.
internal sealed class NullableNodesField : TypeInterceptor
{
    public override void OnBeforeCompleteType(
        ITypeCompletionContext context,
        TypeSystemConfiguration configuration)
    {
        if (configuration is not ObjectTypeConfiguration type
            || type.Name != "Query")
        {
            return;
        }

        foreach (var field in type.Fields)
        {
            if (field.Name is "nodes")
            {
                field.Type = TypeReference.Parse(
                    "[Node]", TypeContext.Output);
            }
        }
    }
}
\end{minted}

\index{Program@\code{Program.cs}!the Speakers registration}%
And that service's \code{Program.cs}, complete, with the same one line marked:

\begin{minted}[highlightlines={65}]{csharp}
using System.Text;
using ChilliCream.Nitro.App;
using HotChocolate.AspNetCore;
using Microsoft.EntityFrameworkCore;
using Microsoft.IdentityModel.Tokens;
using Speakers;

var builder = WebApplication.CreateBuilder(args);

// EF Core reports every command through the logging pipeline as well,
// at Information and with a duration attached. StatementLog below is
// this service's instrument, so silence the other one rather than have
// every statement printed twice.
builder.Logging.AddFilter(
    "Microsoft.EntityFrameworkCore.Database.Command",
    LogLevel.Warning);

builder.Services.AddDbContextFactory<SpeakerContext>(options =>
{
    options.UseSqlite("Data Source=speakers.db");

    if (builder.Environment.IsDevelopment())
    {
        options.AddInterceptors(new StatementLog());
    }
});

// A shared secret rather than a JWKS endpoint, because this book
// stands up no identity provider and a reader has to be able to mint
// a token and repeat every request in the chapter. The three
// Validate flags are off for the same reason and are the three a
// deployment turns on first. What is not negotiable is that the
// router validates the same token against the same key: a subgraph
// that trusts a header the router set, without checking it, has
// moved the guard back out of the service again.
builder.Services
    .AddAuthentication()
    .AddJwtBearer(options =>
    {
        options.TokenValidationParameters = new()
        {
            ValidateIssuer = false,
            ValidateAudience = false,
            ValidateLifetime = false,
            IssuerSigningKey = new SymmetricSecurityKey(
                Encoding.UTF8.GetBytes(SigningKey.Value))
        };
    });

builder.Services.AddAuthorization();

builder
    .AddGraphQL()
    // Without this the [Authorize] attributes compile, reach the
    // schema as @authorize, and are never executed.
    .AddAuthorization()
    .AddSpeakersTypes()
    .RegisterDbContextFactory<SpeakerContext>()
    .AddQueryContext()
    .AddPagingArguments()
    .AddGlobalObjectIdentification()
    .AddApolloFederation()
    .ModifyCostOptions(options => options.ApplyCostDefaults = false)
    .TryAddTypeInterceptor<ShareNodeFields>()
    .TryAddTypeInterceptor<NullableNodesField>();

var app = builder.Build();

using (var scope = app.Services.CreateScope())
{
    var factory = scope.ServiceProvider
        .GetRequiredService<IDbContextFactory<SpeakerContext>>();
    using var db = factory.CreateDbContext();
    db.Database.EnsureCreated();
}

app.UseAuthentication();
app.UseAuthorization();

app.MapGraphQL()
    .WithOptions(options =>
    {
        // Serve the Nitro build that shipped in the package. The
        // default, ServeMode.Latest, downloads the current one
        // from ChilliCream's CDN instead.
        options.Tool.ServeMode = ServeMode.Embedded;
    });

app.RunWithGraphQLCommands(args);
\end{minted}

\index{node@\code{node}!after the change}%
With the interceptor in both services, the mixed list against a stopped
Speakers answers \code{\{"nodes":null\}} and leaves the rest of any response it
appears in alone. That is not a fix for
chapter~\ref{ch:second-third-subgraph}'s finding that \code{node(id:)} cannot
route across a seam, which is a separate problem and still there. It is a fix
for the blast radius when it fails.
\index{node@\code{node}!the plural field's wrapper|)}%
